\section{Setting}\label{sec:setting}

This section fixes what a provider may do, how randomness enters an execution, and
who observes it. The guarantees of \cref{sec:requests} are stated against these
definitions; \cref{sec:excluded} lists what they leave out.

\subsection{Workflows and providers}\label{sec:providers}

The analyses work on a core calculus that keeps the parts of OrchLang that matter
for requests. A \emph{model identity} $m$ stands for everything that determines an
endpoint and its request parameters: provider, model name, output cap $\capm(m)$
and client-side byte cap. A \emph{request expression} $(m, t_0, a_1, t_1, \dots,
a_n, t_n)$ is a prompt template with its arguments, where the $t_j$ are literal text
and each $a_j$ is a literal or a variable. Commands are
\[
c \;::=\; \mathsf{skip} \mid c_1;c_2 \mid x \leftarrow \mathsf{call}\ r \mid
\mathsf{if}\ \pi(x)\ \mathsf{then}\ c_1\ \mathsf{else}\ c_2 \mid \mathsf{retry}\ n\ c
\qquad (n \geq 1),
\]
where a guard $\pi$ is a predicate on strings drawn from the forms the surface
language allows: \code{tokens(x)} compared with a constant, or the truth of a
boolean. A store $\sigma$ maps variables to strings. Variables are partitioned into
public inputs, secrets and local bindings; block scoping is a well-formedness
condition that the type system checks.

\begin{definition}[Provider]\label{def:provider}
A \emph{provider} is a pair $(\Pi, V)$ of functions. For a model identity $m$, a
request text $q$ and a random draw $\omega \in \Omega$, $\Pi(m, q, \omega)$ returns a
response text and the number $o \leq \capm(m)$ of output tokens the provider
reports. For the transcript $T$ of one retry attempt and a draw $\omega$,
$V(T, \omega)$ decides whether the attempt succeeded. We write $\Prov$ for the class
of all providers.
\end{definition}

Nothing else is assumed. $\Pi$ and $V$ may depend on every character of the
request, which is the point: a real model answers what it is asked. A provider
with state, such as a prompt cache or a rate limiter, is included as long as its
state is a function of the requests it has received, since the results below all
establish that two executions send the same requests.

\subsection{Randomness and couplings}\label{sec:couplings}

Randomness is supplied by a \emph{tape} $\rho : K \to \Omega$ of independent draws
indexed by keys. Each random event of an execution (a call, or the decision that
ends a retry attempt) reads the tape at a key that a \emph{keying scheme} computes
from the execution so far. We use three schemes. The \emph{global} scheme keys an
event by its position in the execution; the \emph{model} scheme keys a call by
its model and the number of earlier calls to that model; the \emph{request}
scheme keys a call by its model, its request text and the number of earlier calls
that sent the same text. A scheme is \emph{fresh} if no two events of one
execution read the same key. All three are fresh.

\begin{lemma}[Coupling validity]\label{lem:coupling}
Under a fresh keying scheme, the transcript of an execution driven by a random tape
has the same distribution as the transcript of an execution in which every random
event draws independently from $\Omega$.
\end{lemma}

\begin{proof}
The $i$-th event reads $\rho(k_i)$, where $k_i$ is a function of the events before
it and differs from $k_1, \dots, k_{i-1}$ by freshness. Conditioned on the first
$i-1$ draws, which determine $k_i$, the entry $\rho(k_i)$ has not been read, so it
is independent of them and distributed as $\Omega$. The transcript is the same
function of the sequence of draws in both processes.
\end{proof}

Running two executions on the same tape is therefore a proof device, a coupling in
the sense of probabilistic relational verification~\cite{barthe2009formal,barthe2017coupling}.
It never changes what a single execution does. The mock runtime of
\cref{sec:implementation} implements all three schemes, so that the evaluation can
check each verdict under each coupling.

\subsection{Observers and noninterference}\label{sec:observers}

An execution produces a \emph{history}, the sequence of events
$\mathsf{call}(m, q, \mathit{resp}, o)$ and $\mathsf{attempt}(b)$. An observer is a
function of the history. \Cref{tab:observers} lists the four we consider; each is a
function of the one above it, so security against an observer implies security
against every observer below it. The billed input of a call is
$\tau_m(q) + \mathit{env}_m(q)$ for the model's tokenizer $\tau_m$ and request
envelope $\mathit{env}_m$, both arbitrary functions. A network observer who sees the
sizes of encrypted requests and responses sees a function of the trace.

\begin{table}[t]
\centering
\caption{Observers of an execution. Each row is a function of the row above it.}
\label{tab:observers}
\begin{tabularx}{\textwidth}{@{}lX@{}}
\toprule
Observer & What it sees \\
\midrule
\code{trace} & every event of the history, in order \\
\code{provider} & for each provider, its own events, in order \\
\code{bill} & for each model: number of calls, total input tokens, total output tokens \\
\code{spend} & the sum over models of price times tokens \\
\bottomrule
\end{tabularx}
\end{table}

Two stores are \emph{related}, $\sigma_1 \sim \sigma_2$, when they agree on every
public input, including values that have been declassified.

\begin{definition}[Noninterference]\label{def:ni}
A program $c$ is \emph{$O$-noninterfering} for an observer $O$ if for all related
stores $\sigma_1 \sim \sigma_2$ the distributions of $O(h_1)$ and $O(h_2)$ are equal,
where $h_i$ is the history of $c$ from $\sigma_i$ with independent randomness. It is
\emph{pointwise} $O$-noninterfering under a keying scheme if $O(h_1(\rho)) =
O(h_2(\rho))$ for every tape $\rho$.
\end{definition}

This is the standard two-run formulation of noninterference~\cite{goguen1982security,sabelfeld2003language},
with the observer made explicit. By \cref{lem:coupling}, pointwise noninterference
under any fresh scheme implies noninterference. When a workflow must reveal
something, we measure how much with min-entropy leakage~\cite{smith2009foundations}:
the logarithm of the factor by which an observation increases an adversary's chance
of guessing the secret in one try. The \emph{min-capacity} of a channel is its
largest min-entropy leakage over all prior distributions of the secret, and a
workflow leaks at most $b$ bits to $O$ if the channel from the secrets to $O$'s view
has min-capacity at most $b$.

\subsection{What the guarantee excludes}\label{sec:excluded}

Timing is not part of the history, so it is covered only insofar as it is a
function of the history, as with a cache keyed by requests. Providers whose
randomness is correlated with the secret through a channel outside the program are
excluded, as are observers who see more than the history, such as someone with
access to a provider's internal state. Declassification and endorsement are
trusted: the compiler records the justification the programmer writes but checks
nothing about it, and it does not enforce robustness in the sense of Myers,
Sabelfeld and Zdancewic~\cite{Myers2006Robust}, so untrusted text could in
principle influence what a declassification releases. Declassification is also all
or nothing. A value declassified so that one provider may read it becomes public to
every observer, which is coarser than the setting deserves when the provider is
trusted with content but the network is not (\cref{sec:conclusion}).
